% ---------------------------------------------------------------------------
\section{Pricing Strategy, Limitations and Future Work}
\label{sec:recommendations}
% ---------------------------------------------------------------------------

\subsection{From model output to pricing policy}

The model should be used as a pricing anchor with an explicit range, not as a single
authoritative number. The 80 per cent interval is the natural object to show an
agent: at the median listing it spans roughly a quarter of the estimate, wide enough
to contain a normal negotiation but narrow enough to be informative. The point
estimate should be presented as the centre of that range, and the 90 per cent interval
should drive the escalation decision, because it is the interval whose coverage the
monitoring layer tracks.

The subgroup results imply that a single review threshold is the wrong policy. The
entry band is over-predicted by 4.70 per cent, so a first-time buyer is likely to see
a price above what the property eventually achieves; a conservative adjustment of the
published anchor in that band, or automatic review of entry-band estimates, would
reduce the reputational cost of systematic generosity. The luxury band is
under-predicted by 3.92 per cent with a mean absolute error of 31,238 \USD, more than
three times the entry-band error, so luxury estimates should be routed to a valuer by
default regardless of their nominal confidence. The premium neighbourhood segment
shows almost no bias ($+0.08$ per cent) but a mean absolute error of 23,035 \USD: it
is unbiased on average and unreliable in individual cases, which is exactly the
situation in which a wide interval rather than a point estimate should drive the
conversation.

The confidence classes already implement part of this policy. An estimate whose 90
per cent half-width is at most 18 per cent of the point estimate is high confidence
and may be published directly; between 18 and 30 per cent it is moderate and should
be reviewed for high-value listings; above 30 per cent it is low and must be
reviewed. The segment rules should be layered on top of the class rather than
replacing it, because the two capture different risks: the class captures uncertainty
about the individual property, the segment captures known bias in part of the market.

Two further recommendations follow. The interval width is itself a commercial signal:
a wide interval indicates a heterogeneous micro-market or an unusual property, and it
can be shown to the agent as a prompt to request an appraisal rather than as a defect.
And the model should not be used to advise on renovation: renovation accounts for
6.21 per cent of explanation mass, but the synthetic variables are constructed from
the assessor's remodelling year and the coefficients reflect that construction, so a
claim that a kitchen refurbishment adds a specific percentage to value is not
supported. Finally, the economics should be stated in business units: at the observed
mean asymmetric cost of 16,708 \USD{} per listing, every 1,000 valuations carry
roughly 16.7 million \USD{} of weighted pricing error. That is a historical average
rather than a forecast of savings, but it makes clear that a one-point improvement in
the tolerance hit rate is worth more than a small reduction in mean absolute error.

\subsection{Limitations}

The most important limitation is the sample. With 1,458 transactions and 25
neighbourhoods, the model cannot estimate fine interactions, and the luxury band rests
on 115 sales. Comparing eight families mitigates the risk of a poorly suited
specification, but it cannot create information the sample does not contain.

The second is the synthetic context. The amenity, renovation, marketing and macro
variables encode documented assumptions, and their coefficients measure those
assumptions as much as the Ames market. The realism checks in Section~\ref{sec:data}
show that the generated fields behave plausibly, but plausibility is not validation
against an external source. A production system would replace them with licensed
amenity data, county deed records and a live macro feed, and the model would have to be
re-estimated on those inputs.

The third concerns the error distribution. Log residuals are left-skewed and
heteroscedastic, and the worst five per cent of cases account for 22.1 per cent of
total absolute error, so a single global conformal quantile cannot be correct across
the price range; the entry band under-covers by 8.2 percentage points.

The fourth is temporal. The sample covers a single housing cycle and the model has
never observed a high-interest-rate regime. The replay shows that when the macro
context moves outside the reference window, coverage collapses even though accuracy
holds; after a structural break the model would need recalibration before it could be
trusted, and the retraining policy is the mechanism for doing so.

Two methodological caveats belong here as well. The reported out-of-fold metrics are
computed on the same predictions used to select the model, so they are optimistically
biased by a small amount; the temporal holdout is the independent check and shows an
11 per cent degradation in root mean squared error. The conformal coverage figures
are likewise measured on the calibration residuals, with the monitoring replay as the
independent evidence. Neither caveat undermines the conclusions, but both separate an
honest evaluation from a flattering one.

\subsection{Future work}

Table~\ref{tab:future} lists the extensions the team considers most valuable, in
priority order. The first two address defects in the released system; the remainder
extend its scope. The ranking is by expected business value per unit of effort, not by
methodological novelty.

\begin{table}[htbp]
\centering
\caption{Prioritised future work, with the defect or opportunity it addresses and the
expected benefit.}
\label{tab:future}
\small
\begin{tabularx}{\linewidth}{@{}L{0.042\linewidth}L{0.30\linewidth}L{0.26\linewidth}X@{}}
\toprule
\textbf{No.} & \textbf{Extension} & \textbf{Addresses} & \textbf{Expected benefit} \\
\midrule
1 & Segment-conditional conformal calibration or conformalised quantile regression & Entry and luxury under-coverage & Correct coverage within price bands, not only on average \\
2 & Train on the asymmetric loss (expectile regression or weighted objective) & Loss is reported but not optimised & Align the estimate with the 1.5:1 business cost \\
3 & Public deployment with authentication, rate limiting and health checks & No live URL; open CORS in the demonstration build & Satisfies the deployment criterion; protects the service \\
4 & Replace synthetic context with licensed amenity, deed and macro feeds & Coefficients reflect simulated assumptions & External validity and defensible driver analysis \\
5 & Spatial comparable-sales features and sub-neighbourhood location & Location is represented only by indicators and centroid latitude & Sharper pricing in heterogeneous neighbourhoods \\
6 & Automated retraining with shadow evaluation and promotion gates & Manual retraining decision & Faster recovery after drift, with controlled risk \\
7 & Data-quality monitoring for late, duplicated or malformed listings & The replay assumes a clean feed & Detects operational failures before they reach the model \\
\bottomrule
\end{tabularx}
\end{table}

The first item is the most consequential. The current intervals are honest on average
and wrong in the tails, and the entry band is where the product is most likely to be
used by people least able to absorb a pricing error. Mondrian conformal prediction,
which calibrates within each price band, or conformalised quantile regression, which
adapts interval width to the local error scale, would fix the defect with modest
implementation effort and no change to the point estimator
\citep{romano2019cqr,angelopoulos2023conformal}. The second item is conceptually simple
but requires care: optimising the asymmetric loss directly will push predictions
downward, and the resulting model must be evaluated on coverage and tolerance as well
as on weighted cost, or it will simply learn to under-quote.
